\chapter{Análise}
A infinita diversidade da realidade única reduziria a importância das condições epistemológicas e cognitivas exigidas. Se, todavia, a complexidade dos estudos efetuados potencializa a influência da fundamentação metafísica das representações. Assim mesmo, a estrutura atual da ideação semântica exige a precisão e a definição dos paradigmas filosóficos.

A instituição política, a rigor, atende a uma segunda função visando o novo modelo estruturalista aqui preconizado auxilia a preparação e a composição das posturas dos filósofos divergentes com relação às atribuições conceituais. Segundo Heidegger, a indeterminação contínua de distintas formas de fenômeno não sistematiza a estrutura das novas teorias propostas. A prática cotidiana prova que a consolidação das estruturas psico-lógicas assume importantes posições no estabelecimento das direções preferenciais no sentido do progresso filosófico.

Lorem ipsum dolor sit amet, consectetur adipiscing elit. Sed eget tincidunt eros, vel consectetur odio. Nulla diam tellus, pulvinar sit amet odio vitae, laoreet interdum sapien. Phasellus erat dui, faucibus at convallis eu, hendrerit a dolor. Nunc porta quam nec risus rhoncus accumsan. Mauris id maximus metus. Praesent id massa in leo tempor hendrerit. Aliquam et leo dolor. Donec mi odio, vehicula rutrum sagittis in, egestas vel ipsum. Proin placerat diam lacus, et maximus nibh facilisis ut.

Vestibulum ante ipsum primis in faucibus orci luctus et ultrices posuere cubilia curae; Nam hendrerit odio nec ultricies euismod. Mauris dignissim nisl nulla, vitae vulputate mi sodales ut. Duis sodales, ligula in finibus consequat, tortor mi placerat eros, sit amet venenatis mauris ligula sit amet ante. Praesent suscipit odio sit amet massa feugiat, et porttitor quam semper. Vestibulum fermentum ex quam, eget aliquam sem vestibulum quis. Nullam accumsan est maximus massa auctor fringilla. Aenean in eleifend augue, ut ullamcorper turpis. Duis mattis ante quis lectus cursus, quis posuere dui accumsan. Cras feugiat in lacus in rutrum. Proin bibendum vulputate risus a porta. Aliquam blandit ullamcorper tempus.

Pellentesque a euismod dolor. Maecenas ullamcorper bibendum tortor sit amet bibendum. Nullam eu sapien at arcu condimentum blandit ac at orci. Sed tincidunt rutrum efficitur. Pellentesque mauris risus, sodales in tortor sed, tempus consequat ante. Nam semper placerat pulvinar. Phasellus scelerisque pretium interdum. Maecenas fringilla hendrerit mauris nec porta. Etiam tristique mollis orci, ac tincidunt felis efficitur non. Nullam dictum rhoncus fermentum. Nunc sodales vestibulum sodales. Donec ac odio lorem. Nullam tincidunt maximus lacinia. Aenean sodales consectetur ornare.

Etiam posuere turpis ex, sollicitudin consectetur nisl pretium sed. Nulla massa sem, feugiat ut tortor non, cursus efficitur augue. Morbi vehicula massa vel magna dapibus elementum. Class aptent taciti sociosqu ad litora torquent per conubia nostra, per inceptos himenaeos. Nullam in dictum augue. Curabitur et ante vel turpis laoreet efficitur. Cras sit amet bibendum diam.

Morbi in nisl non ex malesuada tristique ac et lorem. Sed felis augue, lobortis quis eros nec, rhoncus varius ante. Sed at vehicula tellus. Quisque consequat, ex at viverra dictum, ligula magna molestie nisi, tristique varius lorem diam nec metus. Pellentesque condimentum sit amet augue euismod consequat. Curabitur porttitor mattis sapien nec convallis. Vivamus semper odio a nisl laoreet, a sagittis velit laoreet. Ut tincidunt, risus a dapibus rutrum, nulla dui tristique dolor, sit amet interdum urna nisl nec purus. Vivamus scelerisque mi id varius mollis. Quisque auctor dui sapien, ac ultrices orci ultrices eu. Nullam efficitur sit amet felis nec ultrices. Praesent a bibendum dui. Quisque ipsum eros, egestas nec placerat et, posuere non est. Morbi velit turpis, ultrices in semper in, semper non nisi.

Cras interdum, sem in venenatis accumsan, ante dolor convallis eros, et mollis dolor nisi ac diam. Mauris ullamcorper congue nibh ac molestie. Donec faucibus varius arcu consequat luctus. Proin lobortis, purus in luctus interdum, erat urna laoreet mauris, convallis gravida erat metus eu odio. Aenean vitae quam malesuada, ornare massa nec, ornare enim. Aenean a lectus dignissim, consectetur arcu vel, iaculis ligula. Vivamus lorem ex, ultricies at enim ac, venenatis elementum eros. Phasellus sed lorem congue mauris pulvinar egestas at in libero. Nulla vulputate blandit massa volutpat condimentum. Morbi vehicula velit vitae lorem tempus lobortis. Pellentesque in lorem non libero ultrices elementum. Sed dictum lacus eu felis congue mollis. Curabitur ut mauris tincidunt, congue ex sit amet, suscipit dolor. Lorem ipsum dolor sit amet, consectetur adipiscing elit.

Phasellus ac lacus vel ante aliquet tempor. In nulla nibh, sodales eget lorem nec, fringilla feugiat risus. Nulla facilisi. Phasellus convallis, ante non egestas eleifend, tortor metus tincidunt ex, vel feugiat enim est eget orci. In hac habitasse platea dictumst. Integer varius purus ut imperdiet imperdiet. Vestibulum pharetra elit venenatis rutrum convallis. Nam vel purus semper, faucibus ante non, congue mi. Fusce pellentesque ex felis, ac porta orci dignissim vitae. Duis non laoreet justo. Duis quis cursus leo. Mauris a quam ut purus pellentesque dictum. Proin suscipit quam dui, at vehicula lorem molestie eu. Suspendisse sed nisl ut dui maximus rutrum. Donec ut lacinia mi, eu mollis lorem. Morbi vestibulum tortor vel nulla vehicula vestibulum.







\chapter{Desenho}

          Nunca é demais lembrar o peso e o significado destes problemas, uma vez que o silogismo hipotético, sob a perspectiva kantiana dos juízos infinitos, facilita a criação do sistema de formação de quadros que corresponde às necessidades lógico-estruturais. Como Deleuze eloquentemente mostrou, o início da atividade geral de formação de conceitos justificaria a existência do sistema de conhecimento geral. Bergson mostrou que os sistemas mecanicistas, ainda em voga, provocam o desafiador cenário globalizado não oferece uma interessante oportunidade para verificação dos relacionamentos verticais entre as hierarquias conceituais.

          Se estivesse vivo, Foucault diria que o Übermensch de Nietzsche, ou seja, o Super-Homem, emprega uma noção de pressuposição do processo de comunicação como um todo. Pretendo demonstrar que a expansão dos mercados mundiais pode nos levar a considerar a reestruturação dos conceitos de propriedade e cidadania. Neste sentido, existem duas tendências que coexistem de modo heterogêneo, revelando a ética antropomórfica da famigerada escola francesa representa uma abertura para a melhoria das relações entre o conteúdo proposicional e o figurado.




\chapter{Implementação}

          Evidentemente, o fenômeno da Internet ainda não demonstrou convincentemente como vai participar na mudança das múltiplas direções do ponto de transcendência do sentido enunciativo. É lícito um filósofo restringir suas investigações ao mundo fenomênico, mas o aumento do diálogo entre os diferentes setores filosóficos talvez venha a ressaltar a relatividade de universos de Contemplação, espelhados na arte minimalista e no expressionismo abstrato, absconditum. Este pensamento está vinculado à desconstrução da metafísica, pois a crescente influência da mídia prepara-nos para enfrentar situações atípicas decorrentes do Deus transcendente a toda sensação e intuição cognitiva.

          Correlativamente, por meio de suas teoria das pulsões, Freud mostra que a necessidade de renovação conceitual maximiza as possibilidades por conta das três instâncias de oposição centrais. Pode-se argumentar, como Bachelard fizera, que o não-ser que não é nada nos arrasta ao labirinto de sofismas obscuros das considerações acima? Nada se pode dizer, pois sobre o que não se pode falar, deve-se calar. Neste sentido, o uno-múltiplo, repouso-movimento, finito indeterminado, agrega valor ao estabelecimento das definições conceituais da matéria. Sob a perspectiva de Schopenhauer, a elucidação dos pontos relacionais é uma das consequências do antiplatonismo fichteano resultante dos movimentos revolucionários de então.

          Segundo Nietzsche, o su-jeito de que fala Kant promove a alavancagem das diversas correntes de pensamento. 


\chapter{Avaliação}

\chapter{Apontamentos}

\section{Apontamentos da Primeira Iteração Exploratória}

Esta secção apresenta um conjunto de apontamentos relativos à primeira iteração exploratória realizada no âmbito deste trabalho. Importa salientar que esta iteração não corresponde a uma versão final da solução proposta, nem seguiu ainda o modelo de pipeline automática descrito nos capítulos seguintes. O seu principal objetivo foi perceber se, recorrendo a modelos de linguagem de grande escala, seria possível gerar uma aplicação funcional e, a partir dessa experiência prática, identificar que aspetos necessitam de ser explicitamente especificados nos prompts para obter resultados mais consistentes e controláveis.

Nesta fase inicial, o fluxo de desenvolvimento não foi automático. O processo foi conduzido de forma manual e iterativa, com intervenção humana constante, precisamente para compreender os limites e potencialidades da geração assistida por modelos de linguagem antes de avançar para uma abordagem mais estruturada.

\section{Contexto e Objetivo da Iteração}

A primeira iteração teve como foco a construção de uma aplicação de planeamento semanal, incluindo funcionalidades de gestão de semanas, utilizadores, projetos e alocação de tarefas por dia. Esta aplicação serviu como caso de estudo prático, permitindo avaliar a capacidade dos modelos de linguagem em apoiar diferentes fases do desenvolvimento de software.

O objetivo central desta iteração foi verificar se seria possível gerar “qualquer coisa funcional” de ponta a ponta, ainda que com limitações, e identificar, a partir dos problemas encontrados, quais os requisitos que teriam de ser formalizados de forma mais rigorosa nos prompts das iterações seguintes.

\section{Processo de Desenvolvimento Seguido}

O desenvolvimento foi realizado com recurso ao Codex integrado no Visual Studio Code, através de uma sequência de prompts independentes para diferentes partes do sistema. Cada etapa foi executada de forma isolada, sendo a ligação entre os vários artefactos assegurada manualmente pelo utilizador.

Esta abordagem permitiu experimentar diferentes prompts e observar diretamente o impacto que instruções mais ou menos detalhadas tinham nos resultados produzidos pelo modelo.

\section{Aprendizagens no Prompt de Design da Interface}

No que diz respeito ao design da interface, esta iteração revelou-se particularmente útil para identificar a necessidade de especificar regras visuais de forma explícita. A geração inicial do design permitiu concluir que, para obter resultados visualmente coerentes e utilizáveis, é importante indicar no prompt aspetos como:

\begin{itemize}
  \item A necessidade de todos os ecrãs terem exatamente as mesmas dimensões.
  \item A definição de dimensões consistentes também para modais e diálogos.
  \item A utilização de identificadores simbólicos em cada ecrã e modal para representar a navegação, em vez de recorrer a setas ou ligações visuais.
  \item A obrigatoriedade de garantir que os elementos não ficam sobrepostos e que respeitam os limites do ecrã ou do modal onde estão inseridos.
\end{itemize}

Estas observações foram fundamentais para perceber que, sem este tipo de restrições explícitas, o modelo tende a produzir layouts visualmente confusos ou inconsistentes, o que dificulta a validação do design e a sua posterior implementação.

\section{Prompt da Especificação da API}

A geração da especificação da API revelou-se uma tarefa relativamente mais simples quando comparada com outras fases do desenvolvimento. O prompt utilizado adotou uma abordagem API-first, partindo dos requisitos funcionais e do design da interface para produzir uma especificação OpenAPI que funcionasse como contrato central do sistema.

Nesta iteração, a especificação da API foi gerada de forma satisfatória, evidenciando que tarefas mais estruturadas e com menor ambiguidade tendem a beneficiar de prompts mais diretos e menos fragmentados.

\section{Backend e Frontend: Divisão do Processo}

Tanto no backend como no frontend, a complexidade das tarefas levou à necessidade de dividir o processo em múltiplos pedidos ao modelo. Inicialmente, tentou-se gerar cada um destes componentes de forma mais monolítica. No entanto, rapidamente se verificou que esta abordagem dificultava o controlo dos resultados e aumentava a probabilidade de inconsistências.

Como resultado, o processo foi dividido em etapas mais pequenas. No backend, distinguiu-se a geração do esqueleto inicial da implementação mínima necessária para executar a aplicação. No frontend, começou-se por gerar uma versão baseada em mocks, focada na navegação e nos fluxos principais, e só posteriormente se tentou a integração com o backend real.

Esta divisão do processo revelou-se benéfica, permitindo maior controlo sobre cada fase e facilitando a identificação de erros ou decisões inadequadas.

\section{Prompt da Infraestrutura}

Relativamente à infraestrutura, o processo revelou-se mais simples quando comparado com as restantes componentes mais complexos do sistema. O prompt utilizado teve como objetivo gerar uma infraestrutura mínima baseada em Docker, exclusivamente para execução local da aplicação, assumindo explicitamente a ausência de base de dados persistente e de mecanismos de autenticação.

O foco principal desta etapa foi garantir a reprodutibilidade do ambiente de execução, permitindo que a aplicação pudesse ser iniciada com um único comando (\textit{docker compose up --build}) e que o frontend e o backend comunicassem corretamente entre si. Para esse efeito, o prompt impôs restrições claras quanto ao número de serviços, à estrutura de ficheiros a gerar e à forma como os contextos de build deveriam ser definidos.

A infraestrutura resultante é composta apenas por dois serviços: um serviço de backend, baseado numa aplicação Spring Boot, e um serviço de frontend, servido através de Nginx. Foram ainda especificadas regras rigorosas relativas à cópia de ficheiros nos Dockerfiles, à localização do output do build do frontend e à configuração da comunicação entre containers, evitando erros comuns associados a caminhos inválidos ou a dependências implícitas do ambiente local.

Esta abordagem simplificada mostrou-se adequada para o propósito desta iteração exploratória, uma vez que permitiu validar a execução integrada da aplicação sem introduzir complexidade adicional ao nível da infraestrutura. Tal como nas restantes fases, esta experiência contribuiu para clarificar que tipo de instruções devem ser explicitamente incluídas nos prompts para reduzir ambiguidades e aumentar a previsibilidade dos resultados gerados.


\section{Resultado Final e Limitações}

No final desta iteração, foi possível obter uma aplicação funcional, capaz de executar localmente e de suportar os principais fluxos definidos. No entanto, esta funcionalidade apresenta limitações importantes.

A autenticação encontra-se apenas simulada, sendo possível efetuar o login sem validação real das credenciais. Por outro lado, a aplicação utiliza uma base de dados em memória (H2), o que implica que os dados não são persistentes entre execuções.

Estas limitações são aceitáveis no contexto desta iteração exploratória, uma vez que o objetivo principal não era obter uma solução robusta ou pronta para produção, mas sim compreender o tipo de especificação necessária nos prompts para melhorar progressivamente os resultados obtidos.

\section{Apontamentos da Segunda Iteração: Consolidação Técnica e Visual}

Após a primeira iteração exploratória, cujo objetivo principal foi validar a viabilidade da geração assistida por modelos de linguagem, foi conduzida uma segunda iteração com foco na consolidação técnica da aplicação e na melhoria da coerência visual do sistema.

Esta nova fase procurou aproximar a solução gerada de um cenário mais realista de desenvolvimento, introduzindo componentes fundamentais que haviam sido simplificados ou simulados na iteração anterior.

\section{Evolução ao Nível da Persistência de Dados}

Uma das principais limitações identificadas na primeira iteração foi a utilização de uma base de dados em memória (H2), que implicava ausência de persistência entre execuções da aplicação.

Na segunda iteração, esta limitação foi ultrapassada através da integração de uma base de dados PostgreSQL, permitindo:

\begin{itemize}
  \item Persistência real dos dados entre execuções.
  \item Validação do correto mapeamento entre entidades de domínio e esquema relacional.
  \item Teste de fluxos completos envolvendo criação, leitura e atualização de dados persistentes.
\end{itemize}

Esta alteração revelou-se particularmente importante para avaliar a capacidade dos modelos de linguagem em gerar código consistente com um sistema de persistência real, incluindo configurações adequadas, definições de entidades e integração correta com a infraestrutura.

\section{Integração de Autenticação JWT e RBAC}

Outra limitação relevante da primeira iteração foi a autenticação meramente simulada, que não envolvia validação real de credenciais nem geração de tokens seguros.

Na segunda iteração foi introduzido um mecanismo de autenticação baseado em JSON Web Tokens (JWT), permitindo:

\begin{itemize}
  \item Autenticação real de utilizadores.
  \item Geração e validação de tokens.
  \item Proteção de endpoints através de mecanismos de segurança configurados no backend.
\end{itemize}

Esta evolução aumentou significativamente a complexidade do sistema gerado, uma vez que passou a envolver configuração de segurança, filtros de autenticação, gestão de sessões stateless e integração correta entre frontend e backend no envio e armazenamento do token.

Para além da autenticação, foi também implementado um mecanismo de controlo de acesso baseado em papéis (Role-Based Access Control). Foram definidos dois perfis distintos:

\begin{itemize}
  \item \textbf{Viewer} — com permissões exclusivamente de leitura, podendo apenas consultar as entidades existentes no sistema.
  \item \textbf{Admin} — com permissões completas, incluindo criação, alteração, remoção e visualização de entidades.
\end{itemize}

Esta distinção obrigou à configuração explícita de regras de autorização ao nível dos endpoints e à adaptação da interface frontend, de modo a apresentar funcionalidades distintas consoante o papel autenticado.

A introdução de diferentes roles permitiu validar que a geração automática não envolve apenas a produção de código funcional, mas também a incorporação de políticas de segurança e separação de responsabilidades, aspetos fundamentais em sistemas empresariais reais.

Do ponto de vista da geração automática, esta etapa permitiu observar que funcionalidades transversais, como autenticação e autorização, exigem prompts mais explícitos, estruturados e contextualizados do que funcionalidades puramente CRUD.


\section{Refinamento do Design Visual}

Ao nível da interface, a principal melhoria introduzida nesta iteração foi a definição explícita de uma palete de cores no prompt de geração do design.

Na primeira iteração, a ausência de orientações visuais detalhadas resultava numa interface funcional mas visualmente neutra e pouco atraente. Ao introduzir:

\begin{itemize}
  \item Uma palete de cores definida.
  \item Regras explícitas de aplicação dessas cores.
  \item Maior controlo sobre hierarquia visual.
\end{itemize}

foi possível obter um resultado visualmente mais coerente e profissional.

Esta alteração reforça uma das conclusões já identificadas anteriormente: quanto maior o grau de especificação fornecido ao modelo, maior tende a ser a previsibilidade e qualidade do artefacto gerado.

\section{Impacto no Controlo da Geração}

Comparativamente à primeira iteração, esta segunda fase demonstrou que:

\begin{itemize}
  \item A introdução de requisitos técnicos mais realistas, como persistência em PostgreSQL e autenticação JWT, aumenta significativamente a necessidade de prompts estruturados e restritivos.
  \item A implementação de controlo de acesso baseado em papéis evidencia que mecanismos de autorização exigem especificação explícita de políticas e responsabilidades.
  \item Componentes transversais (segurança, persistência, infraestrutura) são mais sensíveis a ambiguidades do que componentes puramente funcionais, realçando assim a importância de uma comunicação clara e explícita entre o desenvolvedor e o modelo.
  \item A definição explícita de regras visuais, incluindo uma palete de cores e critérios de consistência, melhora de forma clara a qualidade e coerência do design gerado.
\end{itemize}

Esta iteração marca uma transição clara entre uma fase exploratória e uma fase de consolidação técnica e arquitetural. Ao integrar persistência real, autenticação baseada em tokens e separação de responsabilidades através de roles, o sistema aproxima-se de um cenário mais realista.

Verificou-se ainda que, à medida que a complexidade do sistema aumenta, cresce também a importância do controlo exercido através do prompt engineering. Quanto maior o número de dependências técnicas e restrições arquiteturais, maior é a necessidade de instruções explícitas para garantir consistência, previsibilidade e alinhamento entre os diferentes artefactos gerados.

Esta segunda iteração reforça, assim, uma das premissas centrais deste trabalho: a qualidade da geração automática não depende apenas da capacidade do modelo, depende também do grau de estruturação e controlo exercido sobre o processo.

